% =============================================================================
% D1 — Caractéristiques du bien informationnel
% Draft v1 — 8 mai 2026
% À intégrer dans Chapter02.tex en remplacement du \wip{} de la section D1
% =============================================================================

\subsection{Caractéristiques du bien informationnel}\label{sec:d1-bien}

\subsubsection{Le champ~: au-delà de la non-rivalité}

\markNEO
La théorie économique du bien informationnel s'est constituée autour d'une propriété unique~: la non-rivalité. Arrow~(1962) montre que l'information, une fois produite, peut être utilisée simultanément par un nombre quelconque d'agents sans dégradation~; Romer~(1990) en tire le fondement de la croissance endogène~; \textcite{jones_nonrivalry_2020} formalisent le cas des données dans un modèle où la non-rivalité est le moteur de rendements croissants à l'échelle. \textcite{shapiro_information_1999} fondent la stratégie des entreprises numériques sur la structure de coûts qu'elle implique~--- coûts fixes élevés, coûts marginaux nuls. \textcite{goldfarb_digital_2019} organisent l'ensemble de l'économie numérique autour de la réduction de cinq types de coûts~--- recherche, réplication, transport, tracking, vérification~--- dont trois reposent sur la non-rivalité. \NEO{Ce cadre traite la non-rivalité comme une propriété inconditionnelle du bien informationnel, donnée une fois pour toutes et indépendante du contexte technique.}

Le chapitre~1 montre que cette caractérisation est incomplète. Deux machines de calcul mécanique de la même époque~--- la calculatrice et la tabulatrice~--- ont des propriétés économiques radicalement différentes~: la première est un bien d'équipement rival, substituable, faiblement complémentaire~; la seconde est un système intégré où la couche informationnelle (les données perforées) est non-rivale mais la couche physique (la machine, les cartes) est rivale, excluable par design et structurellement complémentaire. Un fichier de poids d'un modèle entraîné est non-rival~; l'inférence qui l'active mobilise du compute et de l'électricité en temps réel et est rivale au sens classique. Décrire «~le bien informationnel~» par sa seule non-rivalité, c'est projeter la diversité de ces objets sur un seul axe et perdre les distinctions qui déterminent leurs effets économiques.

Le champ que cette dimension ouvre n'est donc pas la non-rivalité~; c'est l'ensemble des caractéristiques économiques qui distinguent les biens informationnels entre eux et des biens physiques, et la manière dont ces caractéristiques varient selon le type de bien, la technologie et le contexte institutionnel.

\paragraph{Une approche par les caractéristiques}

\markNEO
\textcite{lancaster_new_1966} propose un cadre dans lequel les agents tirent leur utilité non des biens eux-mêmes mais des caractéristiques que ces biens portent~: \NEO{un bien est un vecteur dans un espace à $N$ dimensions, et deux biens de catégories différentes peuvent partager des caractéristiques communes.} L'approche a été développée pour les biens de consommation~; nous l'adaptons ici aux biens informationnels en élargissant les caractéristiques au-delà de la consommation, pour inclure les attributs de production, d'échange et d'accumulation. Cet élargissement est nécessaire parce que ce qui distingue un serveur physique d'un algorithme n'est pas seulement la manière dont ils sont consommés, mais la manière dont ils sont produits, stockés, dépréciés et recombinés.

\textcite{quah_digital_2003} a proposé la taxonomie la plus explicite des biens numériques, identifiant cinq caractéristiques~: non-rivalité, expansibilité infinie, discrétion, aspatialité et recombinabilité. Ce cadre est le plus proche de celui que nous proposons, mais il traite ces caractéristiques comme des propriétés fixes~--- les biens numériques \emph{sont} non-rivaux, expansibles, aspatials. L'observation historique et la diversité des biens informationnels contemporains montrent que ces caractéristiques sont des axes de variation, non des constantes. Le spectre rivalité d'Ostrom~(1990)~--- du bien public pur au bien privé pur, en passant par les communs et les biens de club~--- fournit la bonne intuition~: chaque caractéristique est un continuum, et la position d'un bien sur ce continuum dépend de sa configuration technique et institutionnelle.

Sept caractéristiques sont retenues, issues de quatre traditions~: la théorie des biens publics fournit la \emph{rivalité} et l'\emph{excluabilité} (Samuelson~1954, Ostrom~1990)~; l'économie de la connaissance fournit la \emph{codifiabilité} \parencite{cowan_explicit_2000} et la \emph{recombinabilité} \parencite{weitzman_recombinant_1998,quah_digital_2003}~; l'organisation industrielle fournit la \emph{complémentarité} au sens de la supermodularité \parencite{milgrom_economics_1990}~; la théorie du capital fournit la \emph{dépréciation} \parencite{haskel_capitalism_2018} et la \emph{durabilité} au sens de la distinction bien/service \parencite{hill_goods_1977}. Le critère de sélection est double~: chaque caractéristique est ancrée dans au moins une tradition économique établie, et différencie effectivement les types de biens informationnels entre eux\footnote{L'aspatialité de \textcite{quah_digital_2003} est une caractéristique réelle mais dont la force discriminante est faible au sein des biens informationnels~; elle est traitée en D9 (géographie). Les externalités de réseau \parencite{katz_network_1985} résultent de l'architecture du bien mais produisent leurs effets sur la structure de marché~; elles sont traitées en D2. L'intensité en travail incorporé~--- le ratio travail cristallisé dans le bien / travail nécessaire à son activation, concept d'ascendance marxiste (composition organique du capital)~--- est puissant mais relève de D3 (nature du capital) et D4 (travail) plutôt que des caractéristiques du bien comme output.}.


\subsubsection{Trois mécanismes causaux}

Les sept caractéristiques ne génèrent pas sept mécanismes indépendants~; elles se combinent en tensions structurelles dont chacune produit un mécanisme causal identifiable. Trois tensions se dégagent de l'analyse multi-regard.

\paragraph{La tension de commodification}

\markNEO \markMARX \markINST
La combinaison de la non-rivalité (C1), de l'excluabilité artificielle (C2) et de la durabilité hybride (C6) produit une tension que trois traditions économiques diagnostiquent différemment.

\NEO{La lecture néoclassique identifie une défaillance de marché. Si le coût marginal de reproduction est nul~--- ce que la non-rivalité implique pour la composante stock du bien~--- le prix concurrentiel est nul, et ce prix ne couvre pas les coûts fixes de production.} Arrow~(1962) montre qu'un bien informationnel socialement utile risque de ne pas être produit en régime concurrentiel. La solution canonique est la propriété intellectuelle~: un monopole temporaire qui restaure l'incitation. Mais Arrow montre que ce monopole génère lui-même une sous-production par restriction d'accès~--- un dilemme de second-best sans solution de premier rang dans le cadre standard. \textcite{jones_nonrivalry_2020} formalisent cette tension pour les données~: la non-rivalité implique que les données devraient être partagées (efficacité statique), mais l'absence de droits de propriété détruit l'incitation à les produire (efficacité dynamique).

\MARX{La lecture classique et marxiste reformule ce dilemme en termes de contradiction structurelle.} Si le bien informationnel est non-rival et reproductible à coût nul, sa valeur d'échange tend vers zéro alors que sa valeur d'usage peut être considérable. Maintenir une valeur d'échange pour un tel bien exige un travail permanent de clôture~--- brevets, verrous numériques, interfaces propriétaires, contrôle des données d'entraînement. Ce travail n'est pas un mécanisme d'incitation au sens d'Arrow~; c'est une condition de possibilité de la valorisation marchande, renouvelée à chaque instant contre la tendance naturelle du bien. \textcite{boldrin_against_2008} radicalisent cette lecture en montrant, sur des données historiques longues, que le renforcement des droits de propriété intellectuelle n'est corrélé ni à l'accélération de l'innovation ni à l'augmentation de la production de biens informationnels~: l'exclusion protège des rentes de position, pas un processus d'innovation\footnote{La position de Boldrin et Levine est méthodologiquement contestée~: Moser (2013) trouve des effets positifs de la propriété intellectuelle sur l'innovation dans certains contextes historiques, et Lerner (2009) documente des résultats empiriques mixtes selon les secteurs et les périodes. Le \emph{constat structurel}~--- l'IP n'est pas le mécanisme universel d'incitation que sa justification théorique prétend~--- est plus largement partagé que la conclusion radicale de Boldrin et Levine.}.

\INST{La lecture institutionnaliste ouvre une troisième voie.} L'excluabilité n'est pas une propriété naturelle du bien informationnel mais un choix institutionnel~; dès lors, le binaire marché/État est incomplet. Des régimes de communs gouvernés~--- où l'accès est ouvert mais régulé par des règles collectives~--- constituent une alternative documentée. \textcite{ostrom_governing_1990} et \textcite{hess_understanding_2007} montrent que ces régimes sont viables pour les biens de connaissance. Le logiciel libre, les licences Creative Commons et les modèles \emph{open-weight} en sont des incarnations contemporaines. L'excluabilité se déploie sur un spectre de régimes institutionnels~: propriété privée, club, commun gouverné, bien public~--- et chaque régime génère ses propres mécanismes d'incitation et de coordination.

Le mécanisme se formule ainsi~: la non-rivalité du bien informationnel crée une tension structurelle entre efficience allocative (prix = coût marginal = 0) et incitation à produire (prix $>$ coût moyen). La résolution de cette tension passe par le régime d'excluabilité~--- marché avec IP, commun gouverné, ou hybride~--- et ce régime n'est pas un paramètre technique mais une variable institutionnelle dont la valeur détermine la structure de marché en aval. La condition d'activation est la non-rivalité elle-même~; dès qu'un bien informationnel est rival (un serveur, une bande passante saturée), la tension disparaît et la tarification standard opère. La dualité stock/flux de la caractéristique C6 fait que le même objet peut être simultanément non-rival (en tant que stock copiable) et rival (en tant que flux d'activation), ce qui complique toute tarification uniforme.


\paragraph{La frontière mobile de codification}

\markEVO
La caractéristique de codifiabilité (C3), combinée à la recombinabilité (C7), produit un mécanisme de nature différente~: non plus une tension statique à résoudre, mais une dynamique d'extension.

\textcite{cowan_explicit_2000} montrent que la frontière entre savoir tacite et savoir codifié n'est pas une propriété ontologique mais le résultat d'un processus économique coûteux. \EVO{Transformer un savoir tacite en savoir codifié exige un investissement en formalisation, langage et infrastructure de codage~; cet investissement ne se justifie que si les bénéfices de la codification~--- transmissibilité, stockabilité, recombinabilité~--- excèdent ses coûts.} La frontière se déplace avec la technologie et les incitations~: chaque vague ICT a étendu le domaine du codifiable. Steinwender~(2018) montre que le télégraphe transatlantique a permis la codification en temps réel des spécifications de prix du coton, modifiant la structure des échanges commerciaux~; \textcite{juhasz_spinning_2018} montrent que l'ICT augmente disproportionnellement le commerce de biens intermédiaires parce que leurs spécifications sont plus codifiables.

Le mécanisme est le suivant~: l'extension de la frontière de codification élargit le domaine des biens informationnels et, par conséquent, le domaine où les caractéristiques C1--C2 s'appliquent. Plus on codifie, plus le monde économique est peuplé de biens non-rivaux à excluabilité artificielle~--- avec les tensions que le premier mécanisme identifie. La recombinabilité (C7) amplifie cette dynamique~: chaque nouveau bien codifié augmente le stock recombinable, et le nombre de combinaisons possibles croît de façon combinatoire \parencite{weitzman_recombinant_1998}. La codification et la recombinaison forment une boucle auto-renforçante.

Cowan, David et Foray notent que les savoirs les plus faciles à codifier sont codifiés en premier~; les savoirs résiduels sont par définition les plus coûteux à formaliser. Il y a des rendements décroissants dans la codification~: chaque avancée qui déplace la frontière vers des savoirs plus tacites doit mobiliser des ressources croissantes. Cette observation est purement théorique chez Cowan, David et Foray~; elle acquiert une matérialité frappante dans le cas de l'apprentissage automatique, où l'extension de la codification vers des domaines réputés tacites (diagnostic, rédaction, raisonnement) passe par des entraînements dont le coût croît exponentiellement avec la performance. Ce point croise la dimension~2 (rendements) et la dimension~3 (capital)~; il est posé ici comme conséquence de la caractéristique C3, pas comme résultat de D1.


\paragraph{Le mismatch temporel interne}

\markNEO \markMARX
La combinaison de la complémentarité (C4), de la dépréciation (C5) et de la durabilité (C6) produit un troisième mécanisme~: le mismatch temporel à l'intérieur du même système technique.

\textcite{milgrom_economics_1990} formalisent la supermodularité~: dans un système où la valeur de chaque composant croît avec la présence des autres, l'optimisation séparée ne conduit pas à l'optimum global. Les biens informationnels sont massivement supermodulaires~--- un algorithme sans données est vide, un modèle sans infrastructure d'inférence est dormant, un processeur sans système d'exploitation est inerte. Le degré de complémentarité varie~: un outil logiciel autonome est faiblement complémentaire, un composant d'infrastructure est radicalement systémique. La dimension~3 a montré que ce gradient (isolé/intégré) détermine le degré d'irréversibilité de l'investissement.

\NEO{La théorie néoclassique du capital modélise la dépréciation comme un taux géométrique $\delta$ appliqué au stock.} \textcite{haskel_capitalism_2018} documentent des taux de 25 à 35\,\% par an pour le hardware informatique, contre 5 à 10\,\% pour les structures. \MARX{Marx nomme \emph{dépréciation morale} (\emph{moralische Entwertung}, \emph{Capital}, Livre~II) la perte de valeur d'un bien non par usure physique mais par l'apparition d'un bien supérieur~--- un processus discret et brutal, irréductible à un flux continu.} La dépréciation morale est empiriquement plus juste pour les couches informationnelles rapides (modèles, logiciels) que la dépréciation géométrique~: la sortie d'une nouvelle génération détruit la valeur de la précédente indépendamment de son état physique.

Le mécanisme tient à la coexistence de ces deux modes dans le même système complémentaire. La couche physique (bâtiments, câbles, raccordements) se déprécie par usure~--- lente, continue, prévisible, 15 à 25~ans. La couche informationnelle (logiciels, modèles, algorithmes) se déprécie moralement~--- rapide, discontinue, 12 à 36~mois. Mais les deux couches sont physiquement liées par la complémentarité~: le modèle tourne sur le serveur, le serveur est dans le bâtiment, le bâtiment est raccordé à l'infrastructure. Les décisions d'investissement dans la couche lente engagent sur des décennies~; les décisions dans la couche rapide engagent sur des mois. Ce décalage crée un verrouillage structurel~: les investissements dans les couches lentes contraignent les choix dans les couches rapides, et le renouvellement des couches rapides n'élimine pas l'inertie des couches lentes.

La condition d'activation est le degré de complémentarité. Quand le bien informationnel est faiblement complémentaire (un outil autonome), les couches se déprécient indépendamment et le mismatch est sans conséquence. Quand la complémentarité est forte (un système intégré), les temporalités incompatibles des couches créent une dynamique d'investissement permanente~--- réinvestir dans les couches rapides tout en restant lié aux couches lentes. Ce mécanisme croise la dimension~3 (gradient isolé/intégré) et la dimension~8 (cycles d'investissement en infrastructure).


\subsubsection{Confrontation à l'évidence}

\paragraph{Évidence historique}

Le matériau du chapitre~1 instancie les trois mécanismes de manière différenciée.

La tension de commodification (premier mécanisme) apparaît dès le télégraphe. Le brevet Morse crée un monopole temporaire sur un bien à fort coût fixe et coût marginal faible~; Western Union résout la tension par l'intégration verticale et le contrôle du réseau plutôt que par la tarification au coût marginal. Le logiciel de bureau des années~1980 instancie le même mécanisme~--- coûts de développement élevés, coût de copie nul~--- résolu cette fois par le copyright et les licences. Le mouvement du logiciel libre, à partir des années~1990, instancie la voie d'Ostrom~: gouvernance par des règles collectives plutôt que par la propriété ou le marché. La récurrence du mécanisme à travers des configurations technologiques très différentes (réseau physique, logiciel sur disquette, code en ligne) confirme qu'il est structurel~--- lié aux caractéristiques du bien, pas à la technologie du moment.

La frontière de codification (deuxième mécanisme) se déplace à chaque phase. L'écriture codifie la parole~; le télégraphe codifie les cotations de prix en temps réel \parencite{steinwender_real_2018}~; la mécanographie codifie les procédures administratives~; le traitement de texte codifie la composition typographique. À chaque étape, des savoirs auparavant incorporés dans les agents deviennent des biens informationnels transmissibles et recombinables. La trajectoire est cumulative~: chaque codification ouvre un espace de recombinaison qui accélère la suivante. Le coût de codification varie considérablement~: codifier le prix du coton exige un câble et un opérateur~; codifier le diagnostic médical exige des milliards de paramètres entraînés sur des corpus massifs.

Le mismatch temporel (troisième mécanisme) est documenté dans la dimension~3 pour le cas du data center contemporain, mais le chapitre~1 montre qu'il opère aussi historiquement. Le réseau télégraphique (infrastructure lente, 40~ans de durée de vie) porte des messages dont la valeur est éphémère~; le chemin de fer (infrastructure très lente) est coordonné par un système informationnel (formulaires, rapports, règles) qui se renouvelle en continu. Dans les deux cas, des couches à temporalités incompatibles coexistent dans le même système. Ce qui change avec le numérique contemporain n'est pas la nature du mismatch mais son amplitude~: le rapport entre la durée de vie de l'infrastructure (15--25~ans) et celle des modèles qui y tournent (12--18~mois) est d'un à vingt, contre un à cinq ou un à dix pour les cas historiques.


\paragraph{Évidence contemporaine}

La tension de commodification est aujourd'hui la tension structurante du secteur de l'IA. Le coût de production d'un modèle frontier dépasse 100~millions de dollars\footnote{\textcite{cottier_rapid_2025} estiment le coût d'entraînement des modèles frontier à plusieurs centaines de millions de dollars en 2024--2025.}~; le coût marginal d'une inférence supplémentaire est positif mais faible~; le prix concurrentiel ne couvrirait pas les coûts fixes. Le régime d'excluabilité contemporain est sans précédent~: des modèles \emph{open-weight} (Llama, Mistral, DeepSeek) coexistent à un niveau de performance comparable à celui des modèles propriétaires (GPT-4, Claude) sur de nombreux benchmarks, créant un régime bimodal où la tension de commodification opère simultanément dans les deux directions. Les producteurs de modèles propriétaires investissent massivement dans le travail de clôture (API fermées, données d'entraînement non publiées, intégration verticale dans le compute)~; les producteurs de modèles ouverts financent leur activité par des services complémentaires (hébergement, fine-tuning, support). Ce sont deux résolutions institutionnelles distinctes de la même tension~--- exactement le spectre d'Ostrom.

La codifiabilité fait un saut qualitatif. Les modèles de langage codifient des savoirs réputés tacites~--- diagnostic, rédaction, raisonnement juridique, programmation~--- non par formalisation explicite (comme le faisaient les systèmes experts des années~1980) mais par extraction statistique de patterns à partir de corpus massifs. Le mécanisme de Cowan, David et Foray~(2000) s'applique mais avec une inversion~: le coût de codification n'est plus un coût intellectuel de formalisation mais un coût computationnel d'entraînement, et ce coût croît avec la difficulté du savoir à codifier (les \emph{scaling laws} de \textcite{kaplan_scaling_2020} formalisent cette relation). L'extension de la frontière de codification vers les domaines tacites est simultanément une extension massive du domaine des biens informationnels et une mobilisation massive de ressources computationnelles.

Le mismatch temporel s'observe dans la structure d'investissement des hyperscalers. Les data centers hyperscale ont une durée de vie de 15 à 25~ans~; les GPU qui y sont installés ont une durée de vie économique de 3 à 5~ans~; les modèles qui tournent dessus se déprécient moralement en 12 à 18~mois. Les contrats d'achat d'énergie (PPAs) engagent sur 10 à 15~ans. Quatre temporalités incompatibles~--- bâtiment, processeur, modèle, contrat d'énergie~--- coexistent dans le même système supermodulaire. La littérature sur le capital vintage \parencite{johansen_substitution_1959} traite du putty-clay~; elle ne traite pas du cas où plusieurs couches de putty-clay à rythmes différents sont physiquement inséparables. Ce point croise la dimension~3, qui en fait un développement complet.


\subsubsection{Fait stylisé}

Le bien informationnel n'est pas un objet aux caractéristiques fixes. C'est une famille d'objets caractérisée par un vecteur de sept attributs économiques~--- rivalité, excluabilité, codifiabilité, complémentarité, dépréciation, durabilité, recombinabilité~--- dont la position sur chaque axe varie selon le type de bien, le contexte technologique et le régime institutionnel. Cette variabilité est théoriquement fondée (Lancaster~1966) et empiriquement documentée à travers deux siècles de biens informationnels.

Trois mécanismes structurels se dégagent de l'analyse multi-regard. Premièrement, la tension de commodification~: la non-rivalité du bien informationnel crée un dilemme entre efficience allocative et incitation à produire, résolu par le régime d'excluabilité, qui est une variable institutionnelle et non une propriété technique (Arrow~1962, Boldrin et Levine~2008, Ostrom~1990). Deuxièmement, la frontière mobile de codification~: le domaine des biens informationnels n'est pas donné mais croît avec la technologie, chaque extension élargissant le périmètre où la tension de commodification opère, à un coût croissant (Cowan, David et Foray~2000, Weitzman~1998). Troisièmement, le mismatch temporel interne~: la complémentarité systémique des biens informationnels lie des couches dont les rythmes de dépréciation sont incompatibles~--- usure physique lente pour l'infrastructure, dépréciation morale rapide pour les composants informationnels~--- ce qui crée un verrouillage structurel (Milgrom et Roberts~1990, Marx \emph{Capital}~II, Haskel et Westlake~2018).

Ces trois mécanismes ne sont pas spécifiques à une époque~: ils opèrent du télégraphe au modèle de fondation. Ce qui varie entre les configurations technologiques, c'est leur amplitude~: le ratio coût de première copie / coût de reproduction, le spectre de savoirs codifiables, le rapport entre les temporalités de dépréciation des couches. La question de savoir si certaines configurations du vecteur de caractéristiques produisent des mécanismes qualitativement différents de ceux identifiés par la littérature canonique~--- question que Quah~(2003), Shapiro et Varian~(1999) et Goldfarb et Tucker~(2019) ne posent pas, parce qu'ils traitent les caractéristiques comme fixes~--- relève de l'analyse diachronique et empirique des chapitres adjacents.
